This thesis built and evaluated a local, text-only system for detecting Vietnamese financial phishing. It produced a four-class dataset, two locally trained models, an exported model file that runs on an ordinary laptop, and one evaluation of both models on messages they had never seen. The evidence supports a working prototype. It does not support a claim that the system is ready for real use, and it says nothing about how it would perform on real scam traffic.

\section{Main Findings}

The dataset was built so that no source article can appear in more than one split. Its quality was checked automatically and by hand, and the checks found real problems that were repaired rather than shipped.

Two models were trained locally on the same messages, with one seed. Qwen was trained with 4-bit QLoRA and PhoBERT as a full four-class classifier. On the 220 messages set aside at the start, PhoBERT scored higher overall, at macro F1 0.990892 against Qwen's 0.980493. Neither model produced an unreadable answer, and each made exactly one dangerous mistake, calling a scam safe.

The last of those points is the one that carries the most weight, because on the measure that matters most for a phishing tool the two models tied. Since each was trained only once, the overall gap describes this particular run rather than proving that PhoBERT is reliably better.

The two models also do different jobs. PhoBERT returns a label and nothing else. Qwen returns the label together with a risk level, the phrases it reacted to, and advice on what to do. The comparison covers only the part they share.

\section{Limitations}

The full list can be found in Section~\ref{sec:evidence-limits}, and three of those limits carry more weight than the rest.

The messages are generated from public scam warnings rather than collected from real victims, so nothing here shows how the system behaves on real traffic. Everything was run once with one seed, so there is no way to say how much the numbers would move on a second run. Finally, a review of the reorganised code left six critical findings open, which rules out calling the system production-ready.

\section{Ethical Considerations}

The system classifies and explains suspicious messages; it is not built to write convincing scam messages. Even so, a detection model can be wrong, and an interface that looks confident can make a user feel safer than they should. It should help someone decide, not decide for them.

The dataset was built from public warning articles rather than private conversations, which avoids exposing anyone's messages, but it also limits how realistic the set can be.

The system is tuned to prefer a false alarm over a missed scam, because the costs are not equal: a false alarm wastes a minute, a missed scam can cost someone their savings. That trade-off is easier to judge when errors are reported per class rather than hidden inside one overall score, which is how Chapter~V reports them.

\section{Future Work}
\label{sec:future-work}

Four pieces of work would matter most next.

The first is reviewing the whole dataset by hand rather than a sample. The manual checking done here covered 100 messages and the flagged labels, which was enough to find real problems but not enough to guarantee the rest.

The second is using the two models together rather than choosing between them. PhoBERT classified slightly better on the messages set aside at the start, and Qwen is the only one that can write an explanation, so a system that took the label from PhoBERT and the explanation from Qwen would use each for what it is actually good at. It would cost more memory, and it would need its own evaluation, because nothing measured here shows what happens when two models disagree.

A related option is to replace PhoBERT with ViSoBERT \cite{nguyen2023visobert}, which was trained on Vietnamese social-media text and reported better results than PhoBERT on informal Vietnamese tasks including spam detection. The messages this system handles are informal in exactly that way, as the counts in Section~1.1 show, so it is the more natural starting point than the one used here.

The third is adding OCR so that a user can point the tool at a screenshot instead of copying text out of a message. That is not only a convenience: copying text out of a suspicious message means opening it and selecting text next to a malicious link, and one mistaken tap is all it takes. Reading the text from an image removes that risk. It would need its own evaluation, since OCR on Vietnamese brings its own errors with diacritics and layout.
